\textbf{Context.} Distributed systems built on publish--subscribe middleware such as DDS, ROS~2 or Kafka fail in ways that only show once the assembled system is running. These failures include cascades through shared topics, broker and host outages, and partitions. Engineering teams need a digital counterpart of the \emph{architecture} that can answer what-if questions before deployment and can be checked against the running system afterwards. Such architectural digital twins are still built ad hoc.

\textbf{Objective.} We present an engineering method for an architectural digital twin of publish--subscribe systems. A single typed graph model is the source of every twin artefact: the analysis views, four failure simulators, and an executable container-based testbed that stands in for the physical system. Of the four simulators, the discrete-event queue-flow simulator is the one calibrated against and compared with the testbed. The system is built to the level of a batch digital shadow, and every fidelity claim is bounded to middleware, timing and transport behaviour under declared application behaviour.

\textbf{Method.} The method has four parts.
\begin{enumerate}[nosep,leftmargin=*]
\item A graph metamodel separates structural relations from derived dependencies.
\item A model-to-system transformation generates a Docker and Cyclone DDS deployment from the model. A single fault vocabulary drives both the simulators and container-level chaos actions.
\item A system-to-model path turns OpenTelemetry-shaped traces into calibrated simulator parameters.
\item A fidelity-assessment method compares simulated and observed behaviour in steady state, under faults, and in middleware status. It states explicitly which agreements are independent evidence and which hold by construction.
\end{enumerate}
We place each capability on a digital model, shadow and twin maturity ladder and claim only the levels that are built: a digital model, a generated counterpart and a batch digital shadow.

\textbf{Results.} Model consistency is enforced by executable invariants: the simulators re-derive all seven dependency rules from raw relations and must match the repository edge set.
\begin{itemize}[nosep,leftmargin=*]
\item \emph{Calibration.} On a 74-component scenario, topic rates are recovered within 3\% median error, application service times within 1\%, and inferred trigger policies are never wrong. Under controlled rate and service-time perturbations on three Docker counterparts (including a real-world ROS~2 topology), batch calibration cuts delivery error by 77--96\% and latency error by 10--63\%, with monotonic rolling EMA convergence. Unit round trips on a small brokered system also recover broker service times, $k$-of-$n$ trigger policies and, when the workload is bursty, its burstiness. Under declared periodic workload, neither the staleness window nor burstiness can be recovered from telemetry, and we report both as such.
\item \emph{Smoke test.} On a four-container system, simulated and observed fault responses agree exactly at the stall-set level. The test also exposes a fidelity gap: DDS peer rediscovery after a partition takes about 6\,s, which the simulator does not model.
\item \emph{Fidelity study.} In a pre-registered study on eight emulated systems (2,738 counterpart runs, 10,950 simulator runs), the primary comparison is \emph{inconclusive}: the current simulator and the published one cannot be distinguished in how they rank the impact of publisher faults (Hodges--Lehmann $+0.014$, 95\% CI $[-0.002, 0.094]$). Both, and a closed-form first-order reference, rank at about the counterpart's own repeatability (median $\rho \approx 0.92$ against a ceiling of 0.93). In hindsight this follows from the design: under declared periodic behaviour, a publisher crash removes only its direct readers' deliveries, which is what the reference computes. Descriptively, the current model reproduces broker and host fault rankings (median $\rho = 0.98$) and stall sets (median Jaccard 0.94), and lowers median latency error from 41\% to 25\%.
\end{itemize}

\textbf{Conclusion.} An architectural digital twin can be engineered as a family of model transformations around one metamodel, with fidelity treated as a measured, bounded property rather than an assumption. On publisher ranking, the study could not separate simulation from a closed-form formula. The evidence that the counterpart adds lies in infrastructure faults, latency, perturbed-parameter calibration and middleware behaviour the simulator does not model. The remaining step to a live twin---continuous ingestion, topology-drift repair and production actuation---is designed, not implemented.
